\scanpage{11}
Supposons \deleted{\unsure{tjrs}} que $\deleted{\unsure{K|U}}k$ soit une $\mathbb Q$-algèbre. \deleted{Alors}
Soit $\pi(U)$ \deleted{$\exp\xi$} l'ens des éléments primitifs de $U$.
Alors
\diagram{\deleted{\begin{tikzcd}[ampersand replacement=\&]
\unsure{\pi}(U)\arrow[r] \& \unsure{\underline G|\gamma}
\end{tikzcd}}}
\diagram{\begin{tikzcd}[ampersand replacement=\&]
\pi(U)\arrow[r,shift left,"\exp"] \& \gamma(U)\arrow[l,shift left,"\log"]
\end{tikzcd}}
sont des \uline{bijections} \uline{inverses} \uline{l'une de l'autre}

[N.B. il suffit de vérifier que exp.
applique $\pi(U)$ dans $\gamma(U)$, et log applique
$\gamma(U)$ dans $\pi(U)$]. \uline{Par transport de structure},
\uline{on obtient sur} $\pi(U)$ une loi \uline{de groupe}. Elle est
donnée à l'aide \unsure{des crochets|du crochet|de crochets}, par \uline{une} formule
\uline{universelle} [raison \deleted{\unsure{simple}} à priori] ; \unsure{il y a la} \uline{formule de Campbell-Hausdorff},
% slanted; continuation beginning with bell-Hausdorff in the left margin
\margin{et explicitée dans Séminaire Sophus Lie.}

\textcircled{3}\quad Soit $\unsure{U}$ une $k$-\deleted{hyperalgèbre}\deleted{\unsure{algèbre|bigèbre}}\inserted{algèbre} discrète
(i.e. sans topol. spéc.) \deleted{Alors en général les élts $\xi$
de $I(U)$ ne sont pas nilpotents, et
$\exp\xi$ n'est pas défini. Cependant,} soit
$\Lambda$ une $k$-algèbre ``infinitésimale''.

Alors $\deleted{\unsure{A|\pi}\otimes_k} U\otimes_k\Lambda$ est une $\Lambda$-hyperalgèbre
augmentée ; \deleted{et \unsure{maintenant} les élts
de l'idéal d'augmentation $I(U)\otimes_k\Lambda$
ne sont pas nilpotents, mais ceux de}
% archive sheet number: 10

